\cvsection{Publications \& Talks}

\begin{cventries}

\cventry
{Co-Author}
{Leveraging HPC Infrastructure for AI Workloads – The Role of Kubernetes in AI and HPC}
{Conference}
{May 8, 2024}
{
    \begin{cvitems}
    \item {\url{https://hpckp.org/talks/leveraging-hpc-infrastructure-for-ai-workloads/}}
    \item {Challenges and solutions to integrating AI workloads with HPC infrastructures using Kubernetes.}
    \end{cvitems}
}

\cventry
{Co-Author}
{HEROES - D2.2 Workflows containers}
{Paper}
{Apr 8, 2022}
{
    \begin{cvitems}
    \item {\url{https://heroes-project.eu/wp-content/uploads/2022/07/HEROES-Workflow-Containers.pdf}}
    \item {Description of the approach for the containerization of the software used on the CAE and AI workflows selected as examples for the project.}
    \end{cvitems}
}
\cventry
{Co-Author}
{HEROES - D2.1 CAE and AI Workflows Definition}
{Paper}
{Sep 15, 2021}
{
    \begin{cvitems}
    \item {\url{https://heroes-project.eu/wp-content/uploads/2021/09/HEROES\_CAE-and-AI-Workflows-Definition-Document.pdf}}
    \item {Definition of CAE and AI Workflows selected as examples for the project. AI/HPC Workflow integration: the process of containerizing, optimizing, and integrating new applications and workflows in HEROES.}
    \end{cvitems}
}

\cventry
{Co-Author}
{HEROES - D3.1 Architecture Design}
{Paper}
{Sep 3, 2021}
{
    \begin{cvitems}
    \item {\url{https://heroes-project.eu/wp-content/uploads/2021/09/Architecture-Design-Document.pdf}}
    \item {Description the high-level architecture of the HEROES platform and its main components. The architecture is designed to deliver AI (Artificial Intelligence) and HPC workflows as a Service to end-users and third-party services.}
    \end{cvitems}
}
\end{cventries}
